We begin by specifying some notation and recalling and defining some general tools that play an important role throughout this thesis. We will be working with a wide variety of spaces, many of them having their own inner products and norms. First of all, we denote the Euclidean norm on $\R^m$ by $\left|\ \cdot\ \right|$, for all $m \in \N$. For a more general normed space $X$, where $X$ is different from $\R^m$ for some $m$, we write $\norm{\ \cdot\ }_X$ for its norm. When this norm is induced by an inner product, we denote this inner product by $\inner{\ \cdot\ , \cdot\ }_X$. First we collect a few inequalities that are mainly used in the estimates of our existence and uniqueness proof.
\begin{lemma}{Young's inequality}{}
    For every $a, b \in \R$, we have

    \begin{equation*}
        ab \leq \frac{1}{2} \left(a^2 + b^2\right) .
    \end{equation*}

    \noindent
    More generally, for every $\eps > 0$, we have

    \begin{equation*}
        ab \leq \eps a^2 + \frac{1}{4 \eps} b^2 .
    \end{equation*}
\end{lemma}

\begin{proof}
    Let $a, b \in \R$ and $\eps > 0$. Then

    \begin{equation*}
        0 \leq \left(\sqrt{\eps} a - \frac{1}{2 \sqrt{\eps}} b\right)^2 = \eps a^2 - ab + \frac{1}{4 \eps} b^2 .
    \end{equation*}

    \noindent
    This proves the second inequality. The first inequality is a special case of the second, with $\eps$ set to $\frac{1}{2}$.
\end{proof}

The following is a much more complex inequality, which is used very often in the study of differential equations to derive explicit estimates given an inequality that involves the unknown function.
\todo{move definition absolutely continuous here}

\begin{lemma}{Gronwall's lemma: differential form}{gronwall_dif}
    Let $T > 0$ and let $\eta$ be a non-negative, absolutely continuous function on $[0, T]$, which satisfies
    
    \begin{equation*}
        \eta'(t) \leq \phi(t) \eta(t) + \psi(t)
    \end{equation*}
    
    \noindent
    for almost every $t [0, T]$, where $\phi$ and $\psi$ are non-negative, $L^1$ functions on $[0, T]$. Then

    \begin{equation*}
        \eta(t) \leq e^{\int_{0}^{t} \phi(s) \d s} \left[ \eta(0) + \int_{0}^{t} \psi(s) \d s \right] .
    \end{equation*}

    \noindent
    for all $t \in [0, T]$.
\end{lemma}

This is known as the \emph{differential form} of Gronwall's lemma. A proof can be found in \cite{evans2010partial}. We also have an \emph{integral form}, which is especially useful when dealing with non-differentiable functions. 

\begin{lemma}{Gronwall's lemma: integral form}{gronwall_int}
    Let $T > 0$ and let $y: [0, T] \to \R$ and $a: [0, T] \to \R$ be nonnegative functions such that $a$ and $ay$ belong to $L^1([0, T])$ and

    \begin{equation*}
        y(t) \leq C + \int_{0}^{t} y(s)a(s) \d s
    \end{equation*}

    \noindent
    for some constant $C$ and all $t \in [0, T]$. Then for any $t \in [0, T]$, we have

    \begin{equation*}
        y(t) \leq C e^{\int_{0}^{t} a(s) \d s} .
    \end{equation*}
\end{lemma}

This formulation can be found in \cite{chepyzhov02}.

\begin{definition}{}{}
    Let $\Omega \subseteq \R^d$ be open. A Lebesgue measurable function $f: \Omega \to \R$ is said to be \emph{locally integrable} if

    \begin{equation*}
        \int_K |f| < \infty
    \end{equation*}

    for all compact subsets $K$ of $\Omega$. The space of all locally integrable functions is denoted by $L^1_\text{loc}(\Omega)$, where we identify two functions with each other if they agree almost everywhere in $\Omega$.
\end{definition}
\todo{Put this before FLCV}

An important fact that should be noted is that every function that is square-integrable is also locally integrable. That is, $\Ltwo{\Omega} \subseteq L^{1}_\text{loc}(\Omega)$. For the sake of completeness, we will sometimes state definitions or results with the hypothesis that a function is locally integrable, only to use them on functions that are square-integrable.

We now continue with some concepts that are often used to translate local properties into global results. This method plays a fundamental part in our later discussions about the formulation and analysis of weak forms and solutions to partial differential equations (PDEs).

Let $X$ be a topological space. A function $\phi: X \to \R$ is said to have \emph{compact support} if the set

\begin{equation*}
    \text{supp}(\phi) := \overline{\{ x \in X \mid \phi (x) \neq 0 \}}
\end{equation*}

\noindent
is a compact subset of $X$. For any open set $\Omega \subseteq \R^d$, the space of all infinitely-times continuously differentiable functions with compact support in $\Omega$ is denoted by $C^\infty_c(\Omega)$.

An important property of these functions is that they must be 0 on the boundary of $\Omega$. In the particular context of PDE analysis, they are often referred to as \emph{test functions}. They are used in a wide range of applications, for instance in the field of the Calculus of Variations. A fundamental tool in this theory is the following result, hence its name.

\begin{lemma}{Fundamental Lemma of the Calculus of Variations}{flcv}
    Let $\Omega \subseteq \R^{d}$ be open and $f \in L^{1}_\text{loc}(\Omega)$. If

    \begin{equation*}
        \int_{\Omega} f \phi = 0
    \end{equation*}

    for all $\phi \in C^{\infty}_{c}(\Omega)$, then $f = 0$ almost everywhere in $\Omega$.
\end{lemma}

\begin{proof}
    For every $\eps > 0$, define $\eta_\eps \in C_c^\infty$ as the standard mollifier from \cite[Appendix C.5]{evans2010partial}. Then Theorem 7 from \cite[Appendix C.5]{evans2010partial} implies that

    \begin{equation*}
        \int_\Omega \eta_\eps(x - y) f(y) \d y \to f(x) \qquad \text{as } \eps \to 0
    \end{equation*}

    \noindent
    for almost every $x \in \Omega$. But since for small enough $\eps$ the function $y \mapsto \eta_\eps(x - y)$ is in $C_c^\infty(\Omega)$, this limit is $0$ due to the assumption, i.e. $f(x) = 0$ for almost every $x \in \Omega$.
\end{proof}

Multiplying a PDE with a test function and integrating by parts is sometimes referred to as \emph{testing against} a test function $\phi \in C^\infty_c(\Omega)$. We will see more applications of this in \cref{sec:weak_derivatives}.

When we are differentiating functions in $C^k$ defined in a subset of $\R^d$, we sometimes use the abbreviated notation

\begin{equation*}
    D^\alpha := \frac{\partial^{|\alpha|}}{\partial x_1^{\alpha_1} \dots \partial x_d^{\alpha_d}} ,
\end{equation*}

\noindent
where $\alpha = (\alpha_1, \dots, \alpha_d) \in \N^d$ is called a \emph{multi-index} and $|\alpha| := \alpha_1 + \dots + \alpha_d \leq k$.